\documentclass[12pt,a4paper]{article}
\usepackage[utf8]{inputenc}
\usepackage[T1]{fontenc}
\usepackage[spanish]{babel}
\usepackage{geometry}
\usepackage{hyperref}
\usepackage{longtable}
\usepackage{array}
\usepackage{xcolor}
\usepackage{listings}

\geometry{margin=2.5cm}
\hypersetup{
    colorlinks=true,
    linkcolor=blue,
    urlcolor=blue,
    pdftitle={Documentacion del aplicativo WAF Coraza},
    pdfauthor={Equipo WAF Coraza}
}

\lstset{
  basicstyle=\ttfamily\small,
  breaklines=true,
  frame=single,
  backgroundcolor=\color{gray!8}
}

\title{Guia de Usuario e Instalacion\\WAF Coraza + Dashboard}
\author{Proyecto mi-waf-coraza}
\date{\today}

\begin{document}

\maketitle
\tableofcontents
\newpage

\section{Resumen}
Este aplicativo implementa una \textbf{WAF (Web Application Firewall) con Coraza} en un backend Go y un \textbf{dashboard web} en React para monitoreo y pruebas de ataques.

\textbf{Objetivo principal:} detectar y bloquear trafico malicioso (SQLi, XSS, path traversal, escaneo de rutas administrativas, user-agents ofensivos) y visualizar en interfaz los eventos permitidos y bloqueados en tiempo real.

\section{Enlaces Oficiales del Proyecto}
\begin{itemize}
  \item Backend WAF desplegado (Render): \href{https://waf-coraza.onrender.com}{https://waf-coraza.onrender.com}
  \item Frontend Dashboard (Netlify): \href{https://waf-coraza.netlify.app}{https://waf-coraza.netlify.app}
\end{itemize}

\section{Guia de Usuario}
\subsection{Uso desde navegador}
\begin{enumerate}
  \item Abrir el dashboard: \href{https://waf-coraza.netlify.app}{https://waf-coraza.netlify.app}.
  \item Revisar los paneles de metricas (total, permitidas, bloqueadas).
  \item Ejecutar escenarios de prueba desde el laboratorio WAF.
  \item Validar los resultados en:
  \begin{itemize}
    \item Historial de validacion
    \item Consola de trafico
    \item Actividad reciente
  \end{itemize}
\end{enumerate}

\subsection{Uso desde consola (cualquier equipo)}
El backend esta publicado y puede recibir solicitudes desde cualquier cliente HTTP.

\textbf{Prueba legitima (esperado 200):}
\begin{lstlisting}
curl -i "https://waf-coraza.onrender.com/app"
\end{lstlisting}

\textbf{Pruebas maliciosas (esperado 403):}
\begin{lstlisting}
curl -i "https://waf-coraza.onrender.com/admin"
curl -i "https://waf-coraza.onrender.com/app?q=union%20select%201"
curl -i "https://waf-coraza.onrender.com/app?q=%3Cscript%3Ealert(1)%3C/script%3E"
curl -i "https://waf-coraza.onrender.com/app?file=../../etc/passwd"
\end{lstlisting}

\textbf{Criterio de validacion:}
\begin{itemize}
  \item Solicitudes maliciosas responden HTTP 403.
  \item Solicitudes legitimas responden HTTP 200.
  \item Los eventos se reflejan en el dashboard con estado y reglas asociadas.
\end{itemize}

\section{Arquitectura del Sistema}
\subsection{Estructura del Monorepo}
\begin{itemize}
  \item \texttt{waf/}: backend en Go con middleware Coraza.
  \item \texttt{dashboard/}: frontend en React + TypeScript + Vite.
  \item \texttt{netlify.toml}: configuracion de build y publish para el frontend.
\end{itemize}

\subsection{Diagrama Logico}
\begin{enumerate}
  \item Cliente (navegador o curl) envia solicitud HTTP.
  \item Middleware Coraza procesa conexion, URI, headers y body/ARGS.
  \item Si una regla hace match con accion \texttt{deny}, responde \texttt{403}.
  \item Si no hay interrupcion, la solicitud llega al handler de aplicacion.
  \item Se registra evento en memoria para estadisticas y consola de trafico.
\end{enumerate}

\section{Backend WAF (Go + Coraza)}
\subsection{Tecnologias}
\begin{itemize}
  \item Go (net/http)
  \item Coraza WAF v3
  \item Middleware custom para inspeccion y registro de eventos
\end{itemize}

\subsection{Configuracion clave de Coraza}
La inicializacion de la WAF contempla directivas base de operacion:

\begin{lstlisting}
SecRuleEngine On
SecDefaultAction "phase:2,log,deny,status:403"
SecRequestBodyAccess On
SecRequestBodyLimit 134217728
\end{lstlisting}

\textbf{Puntos clave:}
\begin{itemize}
  \item \texttt{SecRuleEngine On}: habilita ejecucion de reglas.
  \item \texttt{SecDefaultAction}: define accion por defecto de bloqueo y logging en fase 2.
  \item \texttt{SecRequestBodyAccess On}: habilita inspeccion de cuerpo/parametros.
  \item En middleware se ejecuta procesamiento de headers y body para activar reglas \texttt{phase:2} y sobre \texttt{ARGS}.
\end{itemize}

\subsection{Reglas WAF implementadas}
\begin{longtable}{|>{\raggedright\arraybackslash}p{1.5cm}|>{\raggedright\arraybackslash}p{3.5cm}|>{\raggedright\arraybackslash}p{6.2cm}|>{\raggedright\arraybackslash}p{2.5cm}|}
\hline
\textbf{ID} & \textbf{Variable} & \textbf{Patron/Objetivo} & \textbf{Accion} \\
\hline
100 & ARGS & \texttt{@rx (?i)union.*select} (sondeo SQLi) & deny 403 \\
\hline
101 & REQUEST\_HEADERS:User-Agent & \texttt{@rx (?i)(sqlmap|nikto|nmap|masscan)} & deny 403 \\
\hline
102 & REQUEST\_URI & \texttt{@contains /admin} (escaneo admin) & deny 403 \\
\hline
103 & REQUEST\_URI & \texttt{@rx (?i)(\.\./|\.\\\\)} (path traversal) & deny 403 \\
\hline
104 & ARGS & \texttt{@rx (?i)(<script|\%3Cscript)} (XSS) & deny 403 \\
\hline
105 & REQUEST\_URI & \texttt{@contains \%2e\%2e} (traversal encoded) & deny 403 \\
\hline
106 & ARGS & \texttt{@rx (?i)(\.\./|\%2e\%2e\%2f|\%2e\%2e/)} traversal en parametros & deny 403 \\
\hline
\end{longtable}

\subsection{Endpoints del backend}
\begin{itemize}
  \item \texttt{/healthz}: health check del servicio.
  \item \texttt{/app}: endpoint legitimo de aplicacion.
  \item \texttt{/api/stats}: metricas agregadas (total, permitidas, bloqueadas).
  \item \texttt{/api/recent}: eventos recientes en memoria.
\end{itemize}

\subsection{Modelo de evento monitoreado}
Cada solicitud almacena campos como:
\begin{itemize}
  \item metodo, path, user-agent, IP cliente
  \item status HTTP
  \item bandera \texttt{blocked}
  \item IDs y mensajes de reglas coincidentes
  \item transaccion Coraza y tiempo de respuesta (ms)
  \item comando curl reconstruido para consola de trafico
\end{itemize}

\section{Frontend Dashboard (React + TypeScript)}
\subsection{Tecnologias y build}
\begin{itemize}
  \item React + TypeScript + Vite
  \item Variable de entorno principal: \texttt{VITE\_WAF\_API\_URL}
  \item Build y deploy en Netlify
\end{itemize}

\subsection{Puntos clave funcionales del frontend}
\begin{itemize}
  \item Panel de metricas: total, permitidas y \textbf{eventos bloqueados}.
  \item Laboratorio WAF con escenarios predefinidos (control legitimo y ataques).
  \item Historial de validacion con semaforo visual:
  \begin{itemize}
    \item \texttt{403} esperado $\Rightarrow$ \textbf{BLOQUEADO}
    \item \texttt{200/404} en escenario que debe bloquearse $\Rightarrow$ \textbf{NO BLOQUEADO}
    \item escenario legitimo permitido $\Rightarrow$ \textbf{PERMITIDO}
  \end{itemize}
  \item Consola de trafico con comando curl, status y mensaje de regla.
  \item Auto-refresh periodico contra \texttt{/api/stats} y \texttt{/api/recent}.
\end{itemize}

\subsection{Escenarios de prueba del laboratorio}
\begin{enumerate}
  \item Control legitimo: \texttt{/app} (debe permitir).
  \item Escaneo admin: \texttt{/admin} (debe bloquear).
  \item Path traversal: \texttt{/app?file=../../etc/passwd} (debe bloquear).
  \item SQLi probe: query con \texttt{union select} (debe bloquear).
  \item XSS probe: query con \texttt{<script>} (debe bloquear).
\end{enumerate}

\section{Instalacion y Ejecucion (Monorepo)}
\subsection{Requisitos}
\begin{itemize}
  \item Go 1.23 o superior
  \item Node.js 20 o superior
  \item npm 10 o superior
\end{itemize}

\subsection{Clonacion del repositorio}
\begin{lstlisting}
git clone <URL_DEL_REPOSITORIO>
cd waf-coraza
\end{lstlisting}

\subsection{Ejecucion local del backend (waf)}
\begin{lstlisting}
cd waf
go mod tidy
go run .
\end{lstlisting}

Backend local disponible en \texttt{http://localhost:8080}.

\subsection{Ejecucion local del frontend (dashboard)}
En otra terminal:

\begin{lstlisting}
cd dashboard
npm install
\end{lstlisting}

Crear archivo \texttt{dashboard/.env}:
\begin{lstlisting}
VITE_WAF_API_URL=http://localhost:8080
\end{lstlisting}

Iniciar frontend:
\begin{lstlisting}
npm run dev
\end{lstlisting}

Frontend local disponible normalmente en \texttt{http://localhost:5173}.

\section{Operacion y Verificacion}
\subsection{Pruebas recomendadas con curl}
\begin{lstlisting}
# Debe bloquear (403)
curl -i "https://waf-coraza.onrender.com/admin"
curl -i "https://waf-coraza.onrender.com/app?q=union%20select%20user"
curl -i "https://waf-coraza.onrender.com/app?q=%3Cscript%3Ealert(1)%3C/script%3E"
curl -i "https://waf-coraza.onrender.com/app?file=../../etc/passwd"

# Debe permitir (200)
curl -i "https://waf-coraza.onrender.com/app"
\end{lstlisting}

\subsection{Criterio de aceptacion}
Se considera correcto cuando:
\begin{itemize}
  \item Los escenarios maliciosos retornan \texttt{HTTP 403}.
  \item El contador \textbf{Eventos bloqueados} incrementa en dashboard.
  \item En \texttt{Bloqueos recientes} aparecen rule IDs y mensajes de Coraza.
\end{itemize}

\section{Despliegue y Mantenimiento}
\begin{itemize}
  \item Backend productivo en Render: \href{https://waf-coraza.onrender.com}{https://waf-coraza.onrender.com}
  \item Frontend productivo en Netlify: \href{https://waf-coraza.netlify.app}{https://waf-coraza.netlify.app}
  \item Configurar en Netlify la variable \texttt{VITE\_WAF\_API\_URL=https://waf-coraza.onrender.com}.
\end{itemize}

\section{Conclusiones}
El aplicativo implementa una arquitectura completa de defensa y observabilidad:
\begin{itemize}
  \item WAF activa con reglas explicitas para amenazas comunes.
  \item Dashboard para visibilidad operativa y validacion funcional.
  \item Integracion desplegada con enlaces productivos y pruebas reproducibles.
\end{itemize}

\end{document}
